\subsection*{把一个向量变成一把测量尺}
若要回答“一个位移在某方向上走了多少”，输出应该是一个数。固定方向 $u$，把 $v$ 沿它投影就是这种测量。反方向应给负数，垂直方向应给零，相接的两段位移应使测量结果相加。我们因此期待一个线性测量。

在标准正交坐标中，测量由它对基向量的读数决定。若对 $e_1,e_2$ 的读数分别为 $a,b$，对 $(x,y)$ 的读数只能是 $ax+by$。把读数收进向量 $(a,b)$，便得到点积公式。一端是被测对象，另一端是测量仪器，这是对偶空间的具体入口。

\subsection*{区分投影系数、长度与向量}
取 $u=(1,2),v=(3,1)$，点积为 $5$，但投影长度不是 $5$，因为尺子不是单位长度。要找投影向量 $cu$，应要求残差 $v-cu$ 垂直于 $u$。令 $u\cdot(v-cu)=0$，得到 $c=(u\cdot v)/(u\cdot u)$，本例恰为 $1$。

残差 $(2,-1)$ 与尺子点积为零，说明分解正确。这个推导也解释了分母和 $u\ne0$ 的要求。向量的几何长度、投影向量、沿基方向的系数是三个对象，不能因为都出现一个数字就混为一谈。

若改用不垂直的基，逐坐标相乘相加一般不再表示同一几何内积，还需记录基之间的交叉贡献。后面的注意力虽然处理高维数据，仍沿用“用固定方向读出一个分数”的思想；但语义是否与某个方向对应，需要另作经验判断。
